% =====================================================================
%  CONCLUSION AND FUTURE WORK  (unnumbered — HUST convention)
% =====================================================================
\chapter*{CONCLUSION AND FUTURE WORK}
\addcontentsline{toc}{chapter}{CONCLUSION AND FUTURE WORK}

Section~\ref{sec:problem-statement} posed a single question: how can a
physically grounded, end-to-end simulation be constructed to evaluate the
feasibility, signal chain, and array-geometry choices of a future low-cost UHF
solar radio interferometer before hardware deployment? That question was not
academic. It followed directly from my own preliminary two-element
deployment at Xuan~Mai (Section~\ref{sec:motivation}), where phase-incoherent
receivers, cable loss, local RFI, and an under-characterised noise floor
produced an ambiguous result that could not, on its own, be diagnosed. This
thesis answered the question by building and validating a complete software
model of a four-element UHF interferometer at \SI{610}{\mega\hertz}, pursuing
the four objectives set out in Section~\ref{sec:objectives}: antenna
modelling, an end-to-end signal pipeline, a general array-geometry
formulation, and validation against analytic theory and real data. This
chapter reviews what was achieved against each objective, positions the work
relative to existing simulators, restates the contributions, and outlines the
limitations and future extensions that follow.

\section*{1.\quad Summary of Results}

\paragraph{Objective 1 - Antenna modelling.}
The receiving element, a U-frame folded-dipole Yagi-Uda antenna (HKD G11KD-T2), was
modelled in 4NEC2 using the Method of Moments and independently cross-checked using HFSS
(Finite Element Method), directly answering the cross-validation requirement of
Section~\ref{sec:objectives}. The 4NEC2 solution yielded an input impedance of $(108 - j22)\,\si{\ohm}$, a
forward gain of \SI{9.48}{dBi}, a half-power beamwidth of $\approx\SI{52}{\degree}$, and
a standing-wave ratio of $1.55$ at the \SI{75}{\ohm} reference. Direct vector-network-analyser
measurements at the Xuan~Mai test site confirmed the impedance to within calibration
uncertainty, and the HFSS solution agreed with 4NEC2 to better than $0.2~\si{dB}$.
The mutual-coupling study showed that elements are decoupled below $-78~\si{dB}$ at
\SI{5}{\metre} separation and below $-83~\si{dB}$ at the \SI{8}{\metre} minimum array
spacing, formally justifying the independent-sensor assumption used throughout the
pipeline.

\paragraph{Objective 2 - End-to-end signal pipeline.}
An end-to-end Python pipeline was implemented that converts real SWS Learmonth flux
into per-antenna complex baseband voltages through the Friis reception relation~\citep{friis1946,tms2017} and
Johnson-Nyquist thermal noise~\citep{nyquist1928}, applies per-antenna geometric phase derived from the
ENU$\rightarrow$XYZ$\rightarrow\uvw$ transformation~\citep{tms2017}, cross-correlates to form
visibilities, and inverts by a vectorised direct DFT to a dirty image. Two complementary
simulation engines were provided: Engine~A (voltage synthesis, suitable for FPGA export
validation) and Engine~B (forward-model visibilities, analytically exact and
computationally fast for array-geometry studies).

\paragraph{Objective 3 - General array geometry.}
The ENU$\rightarrow$XYZ$\rightarrow\uvw$ formulation was applied identically to every
array configuration evaluated, from the two-element baseline used for fringe validation
to sixteen-element random layouts, without any change to the underlying geometry code.
The parametric sweep exercised this generality directly: five configurations-compact
east-west, north-south line, tee, graded cross, and random-quantified the trade-off
between layout and imaging quality, and a sixth, DLITE-inspired triangle-plus-centre layout
(Case~1.75, Section~\ref{subsec:case175} of Chapter~\ref{ch:pipeline}) showed that a
minimum-redundancy topology recovers two-dimensional resolution from only four elements.
The compact east-west array ($b_{\rm max}=\SI{24}{\metre}$) detects the Sun and
reproduces the correct fringe behaviour but provides only east-west angular resolution.
The extended graded cross ($b_{\rm max}=\SI{130}{\metre}$) achieves two-dimensional
UV coverage sufficient to resolve solar structure and reduces the PSF sidelobe level by
$\approx20~\si{dB}$.

\paragraph{Objective 4 - Validation on 20 March 2025.}
Run on the spring-equinox transit ($\delta=0.0^\circ$, 12-hour window,
$H\in[\pm90^\circ]$), the simulated instrument reproduced three independent analytic
benchmarks: (1) a $\sim$14-minute fringe period on the \SI{8}{\metre} baseline,
matching Eq.~\eqref{eq:fringeperiod} within the sampling resolution; (2) a
$\jinc$-shaped visibility amplitude versus baseline length, including the position of
the first null, confirming that the van Cittert-Zernike relation~\citep{tms2017,taylor1999synthesis} is correctly
implemented; and (3) a dirty image numerically identical (to machine precision) to the
input sky convolved with the independently computed PSF. Agreement on all three
benchmarks simultaneously, using real flux as input, validates the internal
consistency of the pipeline at the analytic level and directly answers the
feasibility question of Section~\ref{sec:problem-statement}: yes, such a simulation
can be built, and it behaves as theory predicts.

\paragraph{Closing the loop on the Xuan~Mai motivation.}
Validating the pipeline at the analytic level is not the same as explaining why the
original field hardware failed, so the simulation was also compared directly against
the two-element analogue trace recorded at Hoa~Lac on the same date
(Section~\ref{sec:real-comparison} of Chapter~\ref{ch:results}). The comparison
attributed the $23.8\times$ gap in measured solar excess to an elevated effective
receiver system temperature ($T_{\rm sys,eff}\approx\SI{3\,600}{\kelvin}$, versus the
\SI{150}{\kelvin} idealised model), and a separate clock-coherence study
(Section~\ref{sec:clock-offset}) showed that a consumer oscillator drift of
$\xi\approx1$~ppm alone would zero the fringe amplitude regardless of receiver noise.
Together these two, independent diagnoses give a quantitative answer to the question
that motivated this thesis: the Xuan~Mai deployment was not undermined by a flawed
concept, but by an uncalibrated low-noise amplifier and the absence of a shared clock
reference-both addressable with modest, identifiable hardware additions.

\section*{2.\quad Position Relative to Existing Simulators}

The most widely used open-source aperture-synthesis simulator, \textsc{APSYNSIM}~\cite{apsynsim},
provides a real-time interactive GUI that is excellent for teaching. This thesis occupies
a complementary but distinct niche, differentiated from \textsc{APSYNSIM} on four axes:

\begin{itemize}[leftmargin=1.5em]
  \item \textbf{Real data.} The pipeline ingests archived solar flux from the Learmonth
        observatory rather than a synthetic Gaussian sky model. The source of truth is
        the Sun itself.
  \item \textbf{Real antenna beam.} The primary beam is the actual 4NEC2 Yagi pattern
        (loaded from a pre-computed \texttt{.npz} file), not the Rayleigh-criterion
        Gaussian. As shown in Chapter~\ref{ch:results}, this difference accounts for
        $\approx25\,\%$ ($\sim\!1~\si{dB}$) more integrated transit power relative to
        the Gaussian approximation.
  \item \textbf{Physics-based noise.} The noise model is grounded in the Johnson-Nyquist
        thermodynamic formula~\citep{nyquist1928}, not an empirical scale factor. The system temperature
        $T_{\rm sys}=\SI{150}{\kelvin}$ and bandwidth $\Delta\nu=\SI{1}{\mega\hertz}$
        are physical quantities that can be measured and that determine the minimum
        detectable flux density from first principles.
  \item \textbf{Reproducibility and hardware path.} The Jupyter notebook constitutes a
        complete, version-controlled, re-runnable record of the simulation. The
        16-bit IQ quantisation stage and LSB-validation check provide a direct path
        to an FPGA hardware correlator-a direction that \textsc{APSYNSIM} does not
        pursue.
\end{itemize}

\section*{3.\quad Contributions}

Matching the five contributions claimed in Section~\ref{sec:contributions-intro}, this
thesis delivers:

\begin{enumerate}[leftmargin=1.8em]
  \item \textbf{A validated 4NEC2 antenna model} of a low-cost commercial UHF
        Yagi-Uda element, cross-checked against direct measurement and an independent
        FEM solution, with a quantitative mutual-coupling study confirming element
        independence across the operational baselines.
  \item \textbf{An end-to-end interferometry simulation engine} that is correct for
        arbitrary two-dimensional array geometries through a single, consistent $\uvw$
        formulation applied identically to per-antenna phase delay, baseline projection,
        and fringe stopping.
  \item \textbf{A real-data-driven validation} demonstrating that the simulated
        instrument, driven by archived Learmonth \SI{610}{\mega\hertz} flux, reproduces
        the classical analytic signatures of aperture synthesis-fringe period, $\jinc$
        visibility, dirty-image consistency-simultaneously and on real data.
  \item \textbf{A structured comparison with \textsc{APSYNSIM}} identifying four concrete
        respects (real data, real antenna beam, physics-based noise, and a reproducible
        hardware-export path) in which this pipeline extends beyond a purely
        educational simulator.
  \item \textbf{A reproducible, single-notebook workflow} that takes a user from an
        archived SWS daily file to a calibrated dirty image in under two minutes of
        computation, with all parameters controlled from a single configuration cell
        and all results exported automatically.
\end{enumerate}

\section*{4.\quad Limitations}

The present implementation makes the following simplifications, each well-justified
for the validation scope of this work but worth noting for future extensions:
the noise model assumes a frequency-flat system temperature and ignores $1/f$ gain
fluctuations; bandwidth smearing is negligible at $\Delta\nu/\nu=1.6\times10^{-3}$
but would become significant for wider channels; ionospheric differential phase is
sub-degree for the baselines used but would need modelling for longer baselines or
solar-storm conditions; no Jones-matrix mutual-coupling correction is applied (the
coupling is $<-78~\si{dB}$, making this contribution negligible); the imaging
stage produces a dirty image rather than a deconvolved sky map; and the voltage
model assumes perfect clock synchronisation ($\tau_k=0$ for all antennas), whereas
real digitisers require a common phase reference to prevent fringe washing
(Section~\ref{sec:clock-offset}). All six limitations point directly to concrete
extensions.

\section*{5.\quad Future Work}

Several extensions follow naturally from this work:

\begin{itemize}[leftmargin=1.5em]
  \item \textbf{FPGA hardware correlator.} The voltage-chain engine already exports
        16-bit IQ samples in hexadecimal format with a validated noise floor of at
        least 20 LSBs. The immediate next step is to stream these samples to an FPGA
        cross-correlator board and compare the hardware fringe amplitude with the
        software result sample-by-sample, providing an end-to-end hardware-in-the-loop
        validation.
  \item \textbf{Physical array upgrade, following the DLITE precedent.} The parametric
        sweep quantified the imaging gain from extending the compact east-west baseline
        set to a graded cross with $b_{\rm max}=\SI{130}{\metre}$, and the Case~1.75
        evaluation (Section~\ref{subsec:case175} of Chapter~\ref{ch:pipeline}) showed
        that a DLITE-inspired~\citep{helmboldt2021dlite} triangle-plus-centre layout
        recovers two-dimensional resolution from exactly four elements, the same count
        used throughout this thesis; this is not a full DLITE system simulation, only
        the minimum-redundancy topology, evaluated at \SI{610}{\mega\hertz} rather than
        DLITE's \SIrange{30}{40}{\mega\hertz} band. Deploying four low-cost LWA-style
        crossed-dipole antennas~\citep{lwa_pricing} at Xuan~Mai or another site in this
        minimum-redundancy arrangement, mirroring the motivation of
        Section~\ref{sec:motivation}, would turn the present detection instrument into a
        genuine solar imager capable of resolving burst sources to a fraction of the
        solar radius.
  \item \textbf{Solar burst recording.} The present pipeline was validated on the
        quiet-Sun continuum. Replacing the quiet-Sun SWS input with a recorded
        Type~II or Type~III burst event would test the pipeline's transient response
        and demonstrate the dynamic-range headroom shown in Table~\ref{tab:detectability}.
  \item \textbf{Multi-frequency operation.} The SWS Learmonth archive covers eight
        spectral channels from \SI{245}{\mega\hertz} to \SI{2695}{\mega\hertz}.
        Extending the pipeline to a multi-channel mode would enable spectro-spatial
        imaging, tracing a burst's frequency drift (and hence electron beam velocity)
        across the corona.
  \item \textbf{\textsc{clean} deconvolution.} Implementing H\"ogbom~\citep{hogbom1974} or
        multi-scale \textsc{clean} as an optional final stage would recover the true
        sky brightness and achieve the noise-limited dynamic range of
        Table~\ref{tab:detectability}, rather than the PSF-sidelobe-limited range of
        the current dirty image.
  \item \textbf{Wider-field imaging and $w$-projection.} For baselines beyond
        $\SI{130}{\metre}$ or field-of-view extensions to tens of degrees, the flat-sky
        approximation breaks down and $w$-projection or faceting must be incorporated.
  \item \textbf{Real-time operation.} The pipeline is currently a post-processing tool.
        Rewriting the correlator stage as a streaming Python process (e.g.\ using
        \texttt{asyncio} or a GPU back-end) would enable real-time image updates
        at the cadence of the SWS one-second data rate.
  \item \textbf{Jones-matrix calibration.} Even though mutual coupling is $<-78~\si{dB}$
        at the operating baselines, incorporating a Jones-matrix element-pattern correction
        would allow the simulator to be used for precision polarimetry experiments at
        shorter baselines.
  \item \textbf{Phase-coherent clock reference.}
        The clock-offset sensitivity study of Chapter~\ref{ch:results}
        (Section~\ref{sec:clock-offset}) showed that a differential oscillator drift of
        $\xi=1$~ppm, typical of a consumer RTL-SDR crystal, produces a coherence
        time $\Delta t_{\rm coh}\approx\SI{1.6}{\milli\second}$ at \SI{610}{\mega\hertz},
        four orders of magnitude shorter than the \SI{43.2}{\second} integration dump.
        Digital cross-correlation between unsynchronised digitisers therefore yields zero
        fringe amplitude regardless of SNR.  The required hardware addition is a
        GPS-disciplined oscillator (GPSDO, $\xi\lesssim1$~ppt) or a shared
        \SI{10}{\mega\hertz} reference signal distributed to all ADC clock inputs.
        The simulation pipeline can extend this study by sweeping
        \texttt{CLOCK\_OFFSET\_NS} and a new \texttt{DRIFT\_PPM} parameter to produce
        a two-dimensional coherence map ($\xi$, $\Delta t_{\rm dump}$) that
        directly specifies the clock-subsystem tolerance budget for the hardware
        correlator design.
\end{itemize}

Returning to the question posed in Section~\ref{sec:problem-statement}: this thesis
shows that a physically grounded, end-to-end simulation can indeed be built to evaluate
the feasibility, signal chain, and array-geometry choices of a low-cost UHF solar radio
interferometer before hardware deployment, and that doing so turns an ambiguous field
failure into a specific, addressable list of hardware requirements. The work delivers a
faithful, reproducible, and real-data-driven model validated from incident flux to final
dirty image, positioned relative to the open-source simulation landscape, and it provides
a solid, well-documented foundation for the four-element DLITE-inspired array upgrade
motivated in Section~\ref{sec:motivation} and for future extensions toward imaging,
burst science, and real-time operation.
\clearpage
